%% aps to author template--please use pdflatex to edit then to pdf--------------
%\documentclass[11pt]{article}
%\documentclass[A4,twoside]{article}
%\usepackage{slashbox}\usepackage{makecell}\usepackage{diagbox}\backslashbox
%\usepackage{aappss}\usepackage{epstopdf}
%\usepackage{luatexja-fontspec}
%\setmainjfont{FandolSong}
%\usepackage{CJK}
%\usepackage{ctex}
%\usepackage{luacode}

\pdfoutput=1
%\documentstyle[11pt,amssymb,epsfig,amsmath,bm]{article}\def\baselinestretch{1.4}
\documentclass[11pt]{article}\def\baselinestretch{1.4}

\parindent 30pt\textheight 9in\topmargin -.7in\textwidth 6.5in
\oddsidemargin 0in\evensidemargin 0in

\usepackage{graphicx}
\usepackage{amsmath}
\usepackage{amssymb}
\usepackage{dcolumn}
\usepackage{bm}
\usepackage{slashed}
\usepackage{color}
\usepackage{authblk}
\usepackage{ulem}
\usepackage{CJK}
%\usepackage{graphicx}
%\usepackage[noadjust]{cite}
%\usepackage[utf8]{inputenc}
%\usepackage{slashed}

\def\k{{\kappa}}
\def\l{{\lambda}}
\def\L{{\Lambda}}
\def\d{{\delta}}
\def\D{{\Delta}}
\def\o{{\omega}}
\def\O{{\Omega}}
\def\e{{\epsilon}}
\def\E{{\Epsilon}}
\def\a{{\alpha}}
\def\b{{\beta}}
\def\c{{\chi}}
\def\f{{\phi}}
\def\g{{\gamma}}
\def\G{{\Gamma}}
\def\j{{\psi}}
\def\vj{{\varphi}}
\def\h{\eta}
\def\p{{\pi}}
\def\P{{\Pi}}
\def\m{{\mu}}
\def\n{{\nu}}
\def\r{{\rho}}
\def\s{{\sigma}}
\def\S{{\Sigma}}
\def\t{{\tau}}
\def\th{{\theta}}
\def\ph{{\phi}}
\def\ps{{\psi}}
\def\Ph{{\Phi}}
\def\Ps{{\Psi}}
\def\x{{\xi}}
\def\P{{\Pi}}
\def\z{{\zeta}}
\def\({\left(}
\def\){\right)}
\def\[{\left[}
\def\]{\right]}

\newcommand{\lag}{\langle}
\newcommand{\rag}{\rangle}
\newcommand{\tb}{{\bar t}}
\newcommand{\pd}{{\partial}}
\newcommand{\dg}{\dagger}
\newcommand{\pr}{\parallel}
\newcommand{\ind}{\text{ind}}
\newcommand{\pb}{\text{PB}}
\newcommand{\tr}{\text{tr}}
\newcommand{\pp}{\perp}
\newcommand{\tf}{\tilde{f}}
\newcommand{\hf}{{\hat f}}
\newcommand{\tS}{\underline{S}}
\newcommand{\sgn}{\text{sgn}}

%\renewcommand\baselinestretch{1.235}\protect
%\renewcommand\baselinestretch{2.0}\protect
%\abovedisplayshortskip 0 pt plus 3pt
%\belowdisplayshortskip 6 pt plus 2pt minus 2pt
%\abovedisplayskip 6 pt plus 2pt minus 2pt
%\belowdisplayskip 6 pt plus 2pt minus 2pt
% \info{2015}{64}{1}{01{}} \infodate{2015.0.0.}{2015.0.0.}
%=================== Text begin here ==============================

\setcounter{footnote}{0}

\date{\today}

\begin{document}

\begin{CJK}{UTF8}{gbsn}

\title{\bf In-medium Electromagnetic Form Factors and Spin Polarizations}
%{\cfundlink}

\author[1]{{Shu Lin}
\thanks{linshu8@mail.sysu.edu.cn}}
\affil[1]{School of Physics and Astronomy, Sun Yat-Sen University, Zhuhai 519082, China}
\author[1]{{Jiayuan Tian}
\thanks{tianjy26@mail2.sysu.edu.cn}}
%\affil[2]{School of Physics and Astronomy, Sun Yat-Sen University, Zhuhai 519082, China}

\maketitle

%\address{3)}{}

\begin{abstract}
We formulate the coupling between fermion spin and background electromagnetic fields using form factors. We show that the vacuum form factors at tree level reproduce the spin polarization effects found in chiral kinetic theory. The vacuum form factors corresponding to spin couplings to perpendicular electric field, parallel and perpendicular magnetic field are degenerate. The degeneracy is expected to be lifted in medium. As an example, we calculate the in-medium QCD radiative correction to the form factors at one-loop order, where we find partial lift of the degeneracy: the spin couplings to parallel and perpendicular magnetic field are different, but the spin couplings to perpendicular electric and parallel magnetic field remain the same. 
\end{abstract}


\newpage

\section{Introduction}

The spin physics of relativistic heavy ion collisions (HIC) has attracted much attention over the past two decades. Based on spin-orbital coupling, early theories predicted that the final particles in off-central heavy ion collisions should be spin polarized \cite{Liang:2004ph,Liang:2004xn}. Indeed, spin polarization has been observed for $\L$ hyperon in low energy collisions \cite{STAR:2017ckg,STAR:2019erd}. Many theoretical efforts have been made in gaining quantitative understanding of the $\L$ hyperon spin polarization from spin-vorticity coupling \cite{Becattini:2013fla, Fang:2016vpj, Li:2017slc, Liu:2019krs, Becattini:2017gcx, Wei:2018zfb, Wu:2019eyi, Fu:2020oxj, Zhang:2019xya, Weickgenannt:2020aaf} and spin-shear coupling \cite{Hidaka:2017auj, Liu:2021uhn, Becattini:2021suc, Fu:2021pok, Becattini:2021iol, Yi:2021ryh, Lin:2022tma}. Spin also couples with electromagnetic fields besides fluid velocity gradient, which are known as spin-magnetic coupling (SBC) and spin Hall effect (SHE) \cite{Liu:2020dxg, Mameda:2022ojk}. The latter has been under intensive investigations in condensed matter systems \cite{RevModPhys.87.1213, PhysRevLett.95.226801}. Despite short lifetime of electromagnetic fields in HIC, effective electric field can also be realized with baryonic density gradient, making SHE potentially observable \cite{Fu:2022oup}.

The interaction of fermions with external electromagnetic fields (EMF) is interesting on its own. In the weak field regime, the spin interaction can be systematically described by spin kinetic theory \cite{Gao:2019znl, Weickgenannt:2019dks, Hattori:2019ahi, Liu:2020flb, Sheng:2020oqs,Yang:2020hri,Weickgenannt:2021cuo,Sheng:2021kfc,Lin:2021mvw}, which is a massive generalization of chiral kinetic theory (CKT), see \cite{Hidaka:2022dmn} for a recent review. In the strong magnetic field regime, the spin interaction can be described by kinetic theory with Landau level basis \cite{Sheng:2017lfu, Lin:2019fqo, Hattori:2016lqx, Fukushima:2019ugr}. The strong magnetic field is also known to have interesting effect in thermodynamics and transport properties \cite{Zhang:2020ben,Fang:2021ndj,Fukushima:2015wck,Hattori:2016lqx,Fukushima:2017lvb,Hattori:2017qih,Hattori:2016njk,Lin:2021sjw,Peng:2023rjj}, see \cite{Hattori:2023egw} for a review and references therein. 

Our present knowledge of spin coupling with EMF is mostly based on CKT and spin kinetic theory, with the SBC and SHE mentioned above dictated by the bare interaction vertex. It is also known that radiative correction can lead to vertex with more structures parameterized by electromagnetic form factors (EMFF). The situation becomes more interesting in medium: the interaction of fermions and EMF is mediated by medium particle, which have a preferred frame. The breaking of Lorentz symmetry leads to vertex with even more structures, which will be parameterized by in-medium form factors. In fact, analogous in-medium gravitational form factors have been discussed and radiative corrections have been considered for spin-vorticity coupling in \cite{Buzzegoli:2021jeh,Lin:2023ass}. We will extend the analysis to EMFF. 

It is worth pointing out a merit of the form factor formulation: while the kinetic description relies on the assumption of weak coupling of fermions with medium, form factor formulation does not have this restriction. It is relevant for phenomenology of spin polarization in HIC, where the produced quark-gluon plasma (QGP) is not a priori guaranteed to be weakly interacting. The form factor formulation offers a mean to describe spin polarization effect at non-perturbative level.

The paper is organized as follows: in Sec.~\ref{sec_2}, we perform field theory calculations of vacuum fermion correlation function at tree level by resumming arbitrary insertions of external EMF. This establishes equivalence with CKT results; in Sec.~\ref{sec_3}, we reformulate the fermion correlation function using vacuum EMFF. We then generalize to the case with radiative correction in medium by defining in-medium EMFF. The analysis gives rise to three structures: spin couplings to perpendicular electric, parallel and perpendicular magnetic field, which are degenerate in vacuum but the degeneracy can be lifted in medium. As an example, we calculate the in-medium QCD radiative correction to EMFF in Sec.~\ref{sec_4}, where we find partial lift of the degeneracy; Sec.~\ref{sec_outlook} is denoted to summary and outlook.


\section{Vacuum fermion correlation function in EMF}\label{sec_2}

In this section, we consider the fermion correlation function in weak external EMF. The weakness of the EMF allows us to perform Wigner transform to obtain the Wigner function
%We consider the Wigner transformed correlation functions for fermions:
\begin{align}
    S^{<(>)}(X=\frac{x_1+x_2}{2}, P)= \int d(x_1-x_2) e^{i P \cdot (x_1-x_2)}S^{<(>)}(x_1,x_2)U(x_2,x_1),
\end{align}
with $\pd_X\ll P$. The external EMF enter through infinite insertions in the lesser(greater) propagators $S^{<(>)}(x_1,x_2)$, as well as through the gauge link $U(x_2,x_1)=e^{-i\int_{x_1}^{x_2} dz^\m A_\m(z)}$ introduced to maintain gauge invariance\footnote{We have absorbed the electromagnetic coupling constant $e$ in the definition of EMF.}. This quantity is extensively studied in CKT by solving from Dirac equation in external EMF \cite{Hidaka:2016yjf, Hidaka:2017auj}. For right-handed fermions, chiral components of lesser propagator $S^{<\mu}$ up to first order in gradient expansion reads:
\begin{align}\label{Sless_CKT}
    S^{<\mu }(X, P)=\frac{1}{2}\tr[\s^\m S^<(X,P)] = -2\pi \left[\delta (P^2)P^\mu f(p_0) + \frac{1}{2}\epsilon^{\mu \nu \alpha \beta }P_\nu F_{\alpha \beta }\delta '(P^2) f(p_0) \right].
\end{align}
The second term corresponds to response to EMF. In particular, the spatial components give
\begin{align}
    S^{<i}=-2\p (p_0B_i-\e^{ijk}P_jE_k)\d'(P^2)f(p_0),
\end{align}
corresponding to SBC and SHE. These effects are not specific to medium, as can be seen from the fact that the corresponding spectral function is independent of the medium. It shall be clear that the medium independence of the spectral function is tied to the vacuum interaction of the fermions and EMF. This medium independence ceases to hold when we consider medium modified interaction in the next section.
%where the second term correspond to the EMF effect. To obtain the above solution, the quasi-particle approximation in CKT is needed, which is only feasible if the coupling in medium is sufficiently weak. Such condition will no longer satisfied in QGP medium. We adopt a field theory analysis approach to transcend the weak coupling limit.
%$\boldsymbol{(by introducing form factor for the fermion-EMF vertex, thus reaching the coupling conditions for QGP medium.)}$

%In field theory, the correction of the correlation function by the external field gradient can be described as the coupling of the EMF to the fermions. We will see in the following computations that by resumming all the tree-level diagrams, the same solutions as CKT can be obtained.
Now we turn to the field theory calculation of $S^<$. It can be calculated directly using the real time formalism, but the diagrammatic analysis is complicated because there are both type $1$ and $2$ couplings between fermions and EMF in Schwinger-Keldysh contour. We will instead calculate $S^{ra}$ using the $ra$-basis, from which we can obtain $S^<$ through the Kubo-Martin-Schwinger (KMS) relation. 
%For simplicity, we focus on the right-handed spinor field in the r/a basis.
The spinor field in $ra$-basis and $12$-basis are related as:
\begin{align}
    \psi _r = \frac{1}{2}(\psi _1 + \psi _2), \quad \psi _a = \psi _1 - \psi _2,
\end{align}
with the $r$ and $a$ fields interpreted as average and fluctuation fields respectively.
%The correlation function in $ra$ basis thus can be given as follow:
%\begin{align}
%    S_{rr} &=\frac{1}{2}(S^< + S^>) \\
%    S_{ar} &=S^F - S^> \\
%    S_{ra} &=S^F - S^< \\
%    S_{aa} &=0
%\end{align}
%Although the above relation is given in coordinate space, it holds after Fourier transforming into phase space. 
In the absence of the EMF (indicated by the superscript $(0)$), the correlation function in thermal equilibrium in the $ra$-basis are given by:
\begin{align}
    S_{rr}^{(0)}(P)&=(\frac{1}{2}+f(P_0))2\p(P\cdot\bar{\s})\d(P^2),\label{S_prop1}\\
    S_{ra}^{(0)}(P)&=\frac{iP \cdot \bar{\sigma}}{P^2+i \eta\,\varepsilon(P_0)}, \label{S_prop2}\\
    S_{ar}^{(0)}(P)&=\frac{iP \cdot \bar{\sigma}}{P^2-i \eta\,\varepsilon(P_0)}.\label{S_prop3}
\end{align}
Without loss of generality, we focus on particles with positive energy thus $\varepsilon(P_0)=+1$. Using the KMS relation, the spectral function is given by:
\begin{align}
    \rho (P) = S_{ra}^{(0)}-S_{ar}^{(0)} = 2\pi (P \cdot \bar{\sigma})\delta (P^2) .\label{spectral_function}
\end{align}
%Starting from the spectral function, all the correlation functions can be derived. Take the lesser propagator for instance, 
The lesser correlation function can then be expressed as
\begin{align}
    S^<_{(0)}(P)=-\tilde{f}(P_0)\rho (P)=-2\pi (P \cdot \bar{\sigma})\delta (P^2)\tilde{f}(P_0) \label{free_lesser_propagator}.
\end{align}

Now we consider the case with EMF. In the $ra$ basis, the external EMF is naturally represented as the average field $A_r$. The correlation function in external EMF can be represented by the sum of a series of Feynman diagrams as shown in Figure~ \ref{tree_level_resum}. Note that for $S^{ra}$ there is only one possible labelings of $ra$, which significantly simplifies the analysis.
\begin{figure}
  \centering
    \includegraphics[width=1\textwidth]{figures/tree_level.png}
  \caption{Retarded fermion correlation function with arbitrary insertions of external electromagnetic fields.}
  %{\bf Better to indicate ra labeling and momentum inflow. The gauge link is not essential.}
   \label{tree_level_resum}
\end{figure}
Note that there are momentum inflow at each vertex (to be denoted by $q$). Since $q\sim\pd_X$ characterizes the gradient, we shall equivalently perform expansion in $q$ in the analysis below. The expansion of the gauge link is simple: taking a straight path for the gauge link as in the CKT approach, we obtain
%where each vertex is given by $-iA\cdot \sigma$ in the total correlation function, and the gradient of background field is performed by its Fourier transform $q$, which is the momentum of background field. In principle, we can obtain all orders of gradient expansion in this way. For simplicity, we consider only the zero and first order of the gradient expansion. Furthermore, the integral path of gauge link is chosen as a straight line from $x_2$ to $x_1$ because of the classical background field:
\begin{align}
    U(x_2,x_1) = e^{-i\int^{x_2}_{x_1}dz \cdot A(z)} = e^{-iA (X)\cdot (x_2-x_1)}+O(\pd_X^2).
\end{align}
%thus the EMF in the gauge link can be regarded as a constant.
The expansion of the remainder of the correlation function requires some efforts: we need to resum diagrams to all order in $A$ but only up to first order in $q$. The result at $O(q^0)$ corresponds to resummation of arbitrary insertions of constant $A$ with no momentum inflow, and the result at $O(q)$ corresponds to the situation with all constant $A$ except one carrying momentum $q$. The detailed calculations can be found in Appendix \ref{sec_app_A}, with the following correlation function at $O(q^0)$ and $O(q)$ respectively
\begin{align}
    O(q^0):&\; S_{ra}^{\m}(K)\equiv\frac{1}{2}\tr[\s^\m S_{ra}^{(A)}(K)] = \frac{iK^\m}{{K}^2+i \eta},\\
    O(q):&\; S_{ra}^{\m}(K)\equiv\frac{1}{2}\tr[\s^\m S_{ra}^{(A)}(K)] = \frac{i\e^{\m\n\r\s}K_\r q_\s A_\n}{({K}^2+i \eta)^2},
\end{align}
with $K$ identified as the gauge invariant kinetic momentum. $S_{ar}$ can be obtained from the above simply by flipping the sign of $i\h$.
%with $K_1=K-\frac{q}{2}$, $K_2=K+\frac{q}{2}$.
%and they turn into even simpler forms, for which we put some of the detailed calculation in \ref{sec_app_A}. The gauge invariant vertex-resumed terms are given as:
%\begin{align}
%    S_{ra}^{(A)}(K) &= \frac{iK'\cdot \bar{\sigma}}{{K'}^2+i \eta}(-iA\cdot \sigma)\frac{iK\cdot \bar{\sigma}}{K^2+i \eta} \label{ra_resummation}\\
%    S_{ar}^{(A)}(K)
%    &= \frac{iK'\cdot \bar{\sigma}}{{K'}^2-i \eta}(-iA\cdot \sigma)\frac{iK\cdot \bar{\sigma}}{K^2-i \eta} \label{ar_resummation}
%\end{align}
%where $K=P-A$ and $K'=P'-A$ is the kinetic momentum of propagator. (\ref{ra_resummation}) and (\ref{ar_resummation}) already contains the gauge link, which change the canonical momentum $P$ into the kinetic momentum $K$. Subtracting (\ref{ra_resummation}) and (\ref{ar_resummation}) leads to
%\begin{align}
%    \rho _{(A)}(K) &= 2\pi \bar{\sigma }^\rho \sigma ^\nu \bar{\sigma }^\sigma (K_\rho K_\sigma +\frac{1}{2}q_\rho K_\sigma +\frac{1}{2}K_\rho q_\sigma)A_\nu \delta '(K^2)
%\end{align}
%\begin{align}
%    S^{<\mu }_{(A)} &= -\frac{1}{2}Tr\left[ \sigma ^\mu\tilde{f}(P_0)\rho (K) \right] \nonumber \\
%    &=2\pi (2K^\mu K^\nu -g^{\mu \nu }K^2-i \epsilon^{\mu \rho \nu \sigma }K_\rho q_\sigma)A_\nu \delta '(K^2)\tilde{f}(K_0) \label{lesser_resummation}
%\end{align}
%The first two terms in (\ref{lesser_resummation}) seem to cause the gauge-dependence for correlation function. They actually come from the redundant calculation for $\mathcal{O}(q^0)$ order during the resummation. By subtracting the redundant terms, we obtain:
%\begin{align}
%    S_{ra}^{(0)}+S_{ra}^{(A)} &=\frac{iK\cdot \bar{\sigma}}{{K}^2+i \eta} + \frac{iK'\cdot \bar{\sigma}}{{K'}^2+i \eta}(-iA \cdot \sigma)\frac{iK\cdot \bar{\sigma}}{K^2+i \eta}
%\end{align}
%Now the gauge dependence is eliminate. Meanwhile, the free propagator $S^{(0)}(P)$ is also corrected to $S^{(0)}(K)$ in the background field. By a similar calculation as (\ref{spectral_function}) and (\ref{free_lesser_propagator}), we obtain a familiar form of lesser propagator:
To obtain the lesser correlation function, we need to use a generalized KMS relation in external EMF. It is complicated by the momentum inflow from the EMF. A detailed account of the KMS relation in external EMF is given in Appendix~\ref{sec_app_B}. For the result at $O(q^0)$, the KMS relation is formally the same as the free theory. 
\begin{align}\label{Sless_q0}
    O(q^0):&\; S^{<\m}(K)=-(S^{ra}-S^{ar})\tilde{f}(K_0)=-2\p K^\m\d(K^2)\tilde{f}(K_0).
\end{align}
For the result at $O(q)$, naive application of the prescription above encounters ambiguity in the choice of energy in the distribution function, which can be $K_0\pm\frac{1}{2}q_0$. However the difference only gives rise to a subleading contribution in the expansion in $q$, thus to the order of our interest the KMS relation is formally the same. Using
\begin{align}
    \frac{1}{(K^2+i\h)^2}-\frac{1}{(K^2-i\h)^2}=2i\p\d'(K^2),
\end{align}
we obtain
\begin{align}\label{Sless_q1}
    O(q):&\; S^{<\m}(K)=-2\p \e^{\m\r\n\s}K_\r q_\s A_\n \d'(K^2)\tilde{f}(K_0).
\end{align}
The sum of \eqref{Sless_q0} and \eqref{Sless_q1} is in perfect agreement with the CKT result \eqref{Sless_CKT} upon the identification $P\to K$.
%We provide a detailed account of the KMS relation in appendix B. In the end, we arrive at the same lesser correlation function as the CKT result
%\begin{align}
%S^{<\mu }_{(0)}+S^{<\mu }_{(A)} = -2\pi \left[ K^\mu \delta (K^2)\tilde{f}(K_0) + \frac{1}{2}\epsilon^{\mu \rho \nu \sigma }K_\rho F_{\nu \sigma }\delta' (K^2)\tilde{f}(K_0) \right].
%\end{align}
%In the above calculation, we've used the KMS relation to build the relation between spectral function and correlation function, which is natural without background field. When a non-zero background field is considered, we need an additional constraint on the 0 component of background field momentum q: $q_0=0$, i.e., fermions do not exchange energy with background EMF(details are showed in Appendix B). Such constraint lead to the fact that every line of propagator in tree level Feynman diagrams has the same distribution function $\tilde{f}(P_0)$. This is actually reasonable: the process by which the correlation function is modified by the background field can be understood as fermions interacting with the background field until the system reaches a steady state, which requires $q_0 << \tau ^{-1}$，where $\tau $ is the relaxation time. Since we only consider the collisionless case, the relaxation time of system is infinite, which result in $q_0 = 0$ \cite{PhysRevD.95.051901}.


\section{In-medium electromagnetic form factors}\label{sec_3}

Let us first reformulate the results in the previous section using EMFF in vacuum. We start with the representation \eqref{O1Oq} reproduced for convenience below
\begin{align}\label{rep}
S_{ra}=\frac{iP_1\cdot\bar{\s}}{P_1^2+i\h}(-i A\cdot\s)\frac{iP_2\cdot\bar{\s}}{P_2^2+i\h},
\end{align}
with $P_1=K+\frac{Q}{2}$ and $P_2=K-\frac{Q}{2}$. We introduce a frame vector $n$ and define the EMF as $E^\m=F^{\m\n}n_\n$ and $B^\m=\frac{1}{2}\e^{\m\n\r\s}n_\n F_{\r\s}$. We shall impose the following restriction to the EMF
\begin{align}
K\cdot E=0\to K\cdot Q \tilde{A}\cdot n-K\cdot \tilde{A} Q\cdot n=0.
\end{align}
This condition is covariant form of $\vec{k}\cdot\vec{E}=0$, which corresponds to that the electric field does no work on the fermions\footnote{This restriction can be lifted if there is dissipation in the system, see \cite{Gagnon:2006hi,Gagnon:2007qt} for dissipative transports in field theory approaches.}. For arbitrary $A$ field, we need to have\footnote{The condition applies to fermion with specific momentum, which is appropriate for retarded correlation function involving no distribution function. For the lesser correlation function involving a distribution, this leads to the vanishing of electric field, which is the Killing condition in equilibrium in CKT.}
\begin{align}\label{Q_cons}
K\cdot Q=0,\quad Q\cdot n=0.
\end{align}
The second condition in \eqref{Q_cons} implies that the resulting EMF are static, while the first condition in \eqref{Q_cons} ensures that $P_{1,2}$ are on-shell up to $O(Q^2)$: $P_{1,2}^2=K^2\pm K\cdot Q+\frac{Q^2}{4}\sim O(Q^2)$. The latter condition allows us to reinterpret the representation \eqref{rep} as a scattering of on-shell fermions off EMF. To see that, we use the spinor representation for right-handed fermions with positive energy
\begin{align}
P_i\cdot\bar{\s}=u(P_i)u(P_i)^\dg,
\end{align}
with $i=1,2$ to rewrite \eqref{rep} as
\begin{align}\label{scat}
\cdots\(u(P_1)^\dg i\s^\m u(P_2)A_\m\)\cdots.
\end{align}
The expression in the bracket is readily identified as a scattering matrix element. So far the frame vector $n$ is arbitrary. Taking $n^\m=(1,0,0,0)$ and using the following explicit relations \cite{Dong:2021fxn},
\begin{align}\label{spinor_products}
&u^\dg(P_1)u(P_2) = \x_1^\dg\x_2(4p_1p_2)^{1/2},\nonumber\\
&u^\dg(P_1)\s^i u(P_2) = \x_1^\dg\x_2(4p_1p_2)^{1/2}\frac{p_1p_{2i}+p_{1i}p_2+i\e^{ijk}p_{2j}p_{1k}}{p_1p_2+\vec{p}_1\cdot\vec{p}_2},
\end{align}
we can obtain a more useful form of the bracket in \eqref{scat} as
\begin{align}
u(P_1)^\dg \(n^\m+\hat{k}^\m+\frac{i\e^{\m\n\r\s} n_\n K_\r Q_\s}{2(K\cdot n)^2}\) u(P_2),
\end{align}
with $\hat{k}=\frac{\vec{k}}{|\vec{k}|}$. The three structures are readily identified as temporal, longitudinal and transverse parts of the vertex. One may check that they satisfy the Ward identity as $Q\cdot n=0$, $Q\cdot\hat{k}=\frac{Q\cdot K}{K\cdot n}-Q\cdot n=0$ and $Q_\m \e^{\m\n\r\s} n_\n K_\r Q_\s=0$.

Now we extend our discussion to in-medium EMFF. Since the medium introduces a preferred frame, it is convenient to use the medium frame vector $u^\m$ to replace the arbitrary frame vector $n^\m$. It turns out the structures obtained above are complete subject to the constraints $K\cdot Q=0$ and $Q\cdot u=0$. Thus the vertex can be parameterized by three EMFF as
\begin{align}%\label{FF_def}
\G^\m=F_0u^\m+F_1\hat{k}^\m+F_2\frac{i\e^{\m\n\r\s}u_\n K_\r Q_\s}{2(K\cdot u)^2} \label{parameterization}
\end{align}
with $F_0$, $F_1$ and $F_2$ being the dimensionless EMFF. It is useful to analyze the transformation of the EMFF under time-reversal symmetry. Since in medium frame $q_0=0$, the EMFF can only a function of $q^2$ and medium parameters etc. Note that $\G^i$ is roughly speaking the spin (as $\G^i$ replaces $\s^i$), which is T-odd. It follows that the EMFF are all T-even, thus we expect the EMFF to be real functions. We shall confirm this by an explicit example \footnote{In the presence of dissipation, the EMFF also depend on relaxation time, which is T-odd, the corresponding EMFF are complex in general.}.
%To simplify the analysis, we restrict ourselves to the case where $q_0=0$, which in coordinate space corresponding to static EMF. As remarked before, the vacuum type interaction does not distinguish EMF with temporal and spatial inhomogeneities, but medium modified interaction does distinguish the two cases.

The physical meaning of the EMFF can be seen by calculating the correlation function. With details collected in Appendix~\ref{sec_app_A}, the effect of the medium modified vertex is the replacement $\s^\m\to\G^\m$ in \eqref{rep}
\begin{align}
S_{ra}=\frac{iP_1\cdot\bar{\s}}{P_1^2+i\h}(-i A\cdot\G)\frac{iP_2\cdot\bar{\s}}{P_2^2+i\h}.
\end{align}
The trace can be evaluated as
\begin{align}
&\frac{1}{2}\tr[P_1\cdot\bar{\s}A\cdot\G P_2\cdot\bar{\s}]\simeq F_2(\vec{B}\cdot\vec{k}),\\ &\frac{1}{2}\tr[\s^i P_1\cdot\bar{\s}A\cdot\G P _2\cdot\bar{\s}]\simeq \(F_0\e^{ijk}E_jk_k+F_1(k_0B^i-\vec{B}\cdot \vec{k}\hat{k}^i)+F_2(\vec{B}\cdot\vec{k})\hat{k}^i\).
\end{align}
We have dropped $O(q^0)$ terms with the same reasoning as in Appendix~\ref{sec_app_A}, that is these $O(q^0)$ terms only contribute to the resummed correlation functions with constant $A$ field. Using the KMS relation as in the previous section (justification of the KMS relation with the modified vertex can be found in Appendix~\ref{sec_app_B}), we easily arrive at
\begin{align}
S^{<0 } &= F_2 \left(\vec{k} \cdot \vec{B}\right) 2\pi \delta '(K^2) f(K_0) ,\label{0_physical_meanings_of_form_factors}\\
S^{<i } &= \left[F_0 \epsilon^{ijk}E_j k_k + F_1 \left( k_0 B^i - \left(\vec{B} \cdot \vec{k}\right)\hat{k}^i\right) + F_2 \left(\vec{B} \cdot \vec{k}\right)\hat{k}^i\right] 2\pi \delta '(K^2) f(K_0). \label{i_physical_meanings_of_form_factors}
\end{align}
%The resummation of the zero-order terms of $q$ will cause the external momentum $p$ to be replaced by kinematic momentum $k$ (redefined as $p$ for convenience). 
In (\ref{0_physical_meanings_of_form_factors}) and (\ref{i_physical_meanings_of_form_factors}), we see that the three form factors give rise to different spin couplings: $F_0$ corresponds to the spin Hall effect; $F_1$ corresponds to the spin-perpendicular magnetic coupling, and $F_2$ corresponds to the spin-parallel magnetic coupling. In vacuum, three form factors read:
%
\begin{align}\label{vac_EMFF}
F_0=1, \quad F_1=1, \quad F_2=1.
\end{align}
%
%This result is not surprising, since in the leading order there doesn't exist a medium to specialize the direction $\hat{p}^i$. 
The fact that all EMFF are identical is of course tied to the vacuum interaction between fermions and EMF. In the next section, we will see medium modified interaction leads to different couplings between fermion spin and the three sources, namely perpendicular electric field, perpendicular/parallel magnetic fields.

\section{QCD radiative corrections to EMFF}\label{sec_4}

%Having obtained the general structures of form factors, we consider the medium modifications to the vacuum form factors in QGP.%, in which the fermion–background field coupling will be radiatively corrected by gluons (represented by thick curves) . 
As an example, we consider modification of the vacuum EMFF in a QGP medium from QCD radiative corrections. Similar to the vacuum case, loop diagrams lead to vertex correction as well as self-energy correction, corresponding to modifications to interaction and incoming/outgoing states respectively. We shall absorb both corrections into modification to the vacuum EMFF \eqref{vac_EMFF}.
%
%There are two part of processes that contribute to the form factors, namely the vertex correction and self-energy correction. In this section we will obtain these two part of form factor respectively.

%We choose the r/a basis as in the previous section.  For vertex correction, this will induced three diagrams with different loop momentum approximation:
We proceed in the $ra$-basis. The one-loop vertex correction comes from three diagrams shown in Figure~\ref{vertex_correction}.
\begin{figure}
  \centering
    \includegraphics[width=.8\textwidth]{figures/vertex_correction.png}
  \caption{One-loop diagrams for vertex correction with three possible labelings in $ra$ (gluons are represented by thick curves)}
  \label{vertex_correction}
\end{figure}
%where the other free propagators in r/a basis are needed：
We need the following propagators in $ra$-basis in addition to \eqref{S_prop1}-\eqref{S_prop3}
\begin{align}
   S_{rr}(P)&=-2\pi (P \cdot \bar{\sigma })\varepsilon(P_0)\left(\tilde{f}(P_0)-\frac{1}{2}\right)\delta (P^2), \\
   D^{\mu \nu }_{rr}(P)&=-2\pi g^{\mu \nu }\varepsilon(P_0)\left(f(P_0)+\frac{1}{2}\right)\delta (P^2), \\
   D^{\mu \nu }_{ra}(P)&=\frac{-ig^{\mu \nu }}{P^2+i \eta \varepsilon(P_0)}, \\
   D^{\mu \nu }_{ar}(P)&=\frac{-ig^{\mu \nu }}{P^2-i \eta \varepsilon(P_0)}, 
\end{align}
where we have chosen the Feynman gauge $(\xi=1)$ for the gluon propagators and suppressed the color factors. %In the high temperature limit, it's generally believed that the contribution of medium is much greater than that of vacuum, that is, 
To further simplify the calculation, we keep only the leading medium contributions to the EMFF by using the hard thermal loop (HTL) approximation. The factor $\frac{1}{2}$ in $S_{rr}$ and $D^{\mu \nu }_{rr}$ will be ignored as the vacuum contribution. Diagram (a) of Figure~ \ref{vertex_correction} is evaluated as:
\begin{align}
   &\int \frac{d^{4} K'}{( 2\pi )^{4}}\left( -ig\sigma ^{\rho } T_{ij}^{a}\right)\frac{i(K'+P_1) \cdot \overline{\sigma }}{(K'+P_1)^{2}-i \varepsilon(K'_0+P_1^0)\eta }\left( -i\sigma ^{\nu } A_{\nu }\right)\frac{i(K'+P_2) \cdot \overline{\sigma }}{(K'+P_2)^{2}-i \varepsilon(K'_0+P_2^0)\eta }\left( -ig\sigma ^{\sigma } T_{ij}^{a}\right)\nonumber\\
   &\times( -g_{\rho \sigma }) 2\pi \delta \left( K'^{2}\right) f( K'_{0}).
\end{align}
In the HTL approximation, we ignore $P_2$ and $P_1$ next to $K'$ in the numerators. The denominators can be simplified using $\d(K'^2)$ and on-shell conditions of $P_2$ and $P_1$. We then have
%The imaginary part of denominator $i \eta $ actually contributes to two kinds of IR divergence. For simplicity, we can ignore $i \eta $, and interpret the singularity as the collinear divergence: 
\begin{align}\label{diag_a}
   -2ig^2C_F A_\nu \int\frac{d^{4} K'}{( 2\pi )^{4}} 2K'^\nu (K' \cdot \sigma) \frac{1}{2K'\cdot P_1-i \varepsilon(K'_0)\eta}\frac{1}{2K'\cdot P_2-i \varepsilon(K'_0)\eta}2\pi \delta (K'^2)f(K'_0).
\end{align}
There is potential divergence when the momenta $K'$ and $P_1$($P_2$) become collinear. The collinear divergence does not contribute to EMFF as we shall see below, allowing us to ignore $i\h$ here. Similarly diagrams (b) and (c) are evaluated as
%where the factor $\frac{1}{2K\cdot P'}$ and $\frac{1}{2K\cdot P}$ contribute to a logarithmic divergence respectively when $K \parallel P$ or $K \parallel P'$, and a higher order divergence when $K \parallel P \&\& K \parallel P'$.  Similarly, (b) and (c) in Figure \ref{vertex_correction}. can be obtained:
\begin{align}\label{diag_b}
   (b)&:\int \frac{d^{4} K'}{( 2\pi )^{4}}\left( -ig\sigma ^{\rho } T_{ij}^{a}\right)K' \cdot \overline{\sigma}\left( -i\sigma ^{\nu } A_{\nu }\right)\frac{i(K'-Q) \cdot \overline{\sigma }}{(K'-Q)^{2}+i \eta }\left( -ig\sigma ^{\sigma } T_{ij}^{a}\right)\frac{-ig_{\rho \sigma }}{(K'+P_1)^{2}-i \eta } 2\pi \delta \left( K'^{2}\right) \tilde{f}( K'_{0}) \nonumber\\
   &\simeq 2ig^2C_F A_\nu \int\frac{d^{4} K'}{( 2\pi )^{4}} 2K'^\nu (K' \cdot \sigma) \frac{1}{2K'\cdot Q}\frac{1}{2K'\cdot P_1}2\pi \delta (K'^2)\tilde{f}(K'_0),
\end{align}
and
\begin{align}\label{diag_c}
   (c)&:\int \frac{d^{4} K'}{( 2\pi )^{4}}\left( -ig\sigma ^{\rho } T_{ij}^{a}\right)\frac{i(K'+Q) \cdot \overline{\sigma }}{(K'+Q)^{2}-i \eta }\left( -i\sigma ^{\nu } A_{\nu }\right) K' \cdot \overline{\sigma} \left( -ig\sigma ^{\sigma } T_{ij}^{a}\right)\frac{-ig_{\rho \sigma }}{(K'-P_2)^{2}-i \eta } 2\pi \delta \left( K'^{2}\right) \tilde{f}( K'_{0}) \nonumber\\
   &\simeq 2ig^2C_F A_\nu \int\frac{d^{4} K'}{( 2\pi )^{4}} 2K'^\nu (K' \cdot \sigma) \frac{1}{2K'\cdot Q}\frac{1}{2K'\cdot P_2}2\pi \delta (K'^2)\tilde{f}(K'_0).
\end{align}
Here $Q^2$ are dropped in the denominators as $O(q^2)$ subleading corrections.
Summing over \eqref{diag_a}, \eqref{diag_b} and \eqref{diag_c}, we obtain
\begin{align}\label{diag_abc}
   &-2ig^2C_F A_\nu \int\frac{d^{4} K'}{( 2\pi )^{4}} 2K'^\nu (K' \cdot \sigma) \left[ \frac{1}{2K'\cdot P_1}\frac{1}{2K'\cdot P_2}\(f(K'_0)+\tilde{f}(K'_0)\) \right]2\pi \delta (K'^2) \nonumber\\
   &=-2ig^2C_F A_\nu \int\frac{d^{4} K'}{( 2\pi )^{4}} 2K'^\nu (K' \cdot \sigma)\left(\frac{1}{2K'\cdot P_2}-\frac{1}{2K'\cdot P_1}\right)\frac{1}{2K'\cdot Q}\left(f(K'_0)+\tilde{f}(K'_0)\right)2\pi \delta (K'^2). %\label{sum_over_three}
\end{align}
%Note that at the linear order of $q$, the external momentum $P$ and $P'$ can be put on shell. The dirac structure of external propagator thus can be decomposed as $P\cdot \overline{\sigma }=u(P)u^\dagger(P)$. This means that the equation of motion (EOM) is available. We can introduce an extra term for $K^\nu $ in the integrand which is not contribute to the form factor but is able to cancel the logarithmic divergence\cite{Lin:2023ass}:
Let us discuss the collinear divergences using the second line of \eqref{diag_abc}. The two terms in the second bracket of \eqref{diag_abc} are simply related by $P_1\leftrightarrow P_2$. In fact the collinear divergence does not contribute to EMFF for the following reason: since the vertex function is sandwiched between incoming/outgoing right-handed fermion states as $u^\dg(P_1)\G^\m u(P_2)$, terms like $P_2\cdot \s$ and $P_1\cdot \s$ simply does not contribute to the amplitude by equation of motion (EOM). So we may subtract appropriate structures to eliminate the collinear divergence as
\begin{align}
&-2ig^2C_F A_\nu \int\frac{d^{4} K'}{( 2\pi )^{4}} 2 K'^\nu \left(K'^\lambda -\frac{K' \cdot u}{P_2 \cdot u}P_2^\lambda \right) \sigma_\lambda
\frac{1}{2K'\cdot P_2}\frac{1}{2K'\cdot Q}\left(f(K'_0)+\tilde{f}(K'_0)\right)2\pi \delta (K'^2)\nonumber\\
&-(P_2\to P_1) ,\label{integral_with_counterterm}
\end{align}
where $u^\mu =(1, 0, 0, 0)$ denotes the QGP frame vector. The subtracted structure ensures that $K'^\lambda -\frac{K' \cdot u}{P_{i} \cdot u}P_{i}^\lambda$ and $2K'\cdot P_{i}$ approach zero at the same rate, making the integrand regular in the collinear limit. There is actually another divergence occurring when $q\to0$ (note that $q_0=0$ from our restriction on the EMF). This is an IR divergence, which will be cut off by screening effect of the medium.

To perform the tensor integrals, we adopt the following parameterization of the momenta by using the rotational symmetry:
\begin{align}
   P_1 &=(p_0, 0, \frac{q}{2}, p) ,\nonumber \\
   P_2 &=(p_0, 0, -\frac{q}{2}, p) ,\nonumber \\
   Q &=(0, 0, q, 0) ,\nonumber \\
   K' &=(k'_0, k' \sin{\theta } \cos{\varphi }, k' \sin{\theta } \sin{\varphi }, k' \cos{\theta }).\label{parameterization}
\end{align}
One of the advantages of this parameterization is that many of the components in the integral can be eliminated, leaving only those which have definite physical meaning.
After some algebra, we can see that most of the components in (\ref{integral_with_counterterm}) are zero after integration\footnote{The components with $\lambda = 2$ can be written as $Q\cdot \sigma$ in the parameterization (\ref{parameterization}), which also have no contribution to the EMFF by the EOM.  }, leaving only $(\lambda, \nu)=(1, 1), (3, 3), (3, 0)$: 
\begin{align}
(\ref{integral_with_counterterm}) = \left\{
\begin{aligned}
-2i A_\nu \frac{m_f^2}{4p^2}\ln{\frac{2p}{q}}, \quad\quad\quad\quad   & (\lambda,\nu)=(1,1)  \\
-2i A_\nu \frac{m_f^2}{4p^2}\left(1-\ln{\frac{2p}{q}}\right), \quad\quad   & (\lambda,\nu)=(3,3)  \\
-2i A_\nu \frac{m_f^2}{4p^2}\left(-\ln{\frac{2p}{q}}\right), \quad\quad   & (\lambda,\nu)=(3,0)
\end{aligned}
\right.
\end{align}
with
\begin{align}
m_f^2=g^2C_F\int\frac{4\p kdk}{(2\pi)^3}\left(f(K_0)+\tilde{f}(K_0)\right)=\frac{1}{8}g^2 T^2 C_F,
\end{align}
is the thermal mass of fermions with $C_F=\frac{4}{3}$.
Defining $\hat{l}^i=\frac{1}{pq}\epsilon^{ijk}q_j p_k$, we can rewrite the integral in (\ref{integral_with_counterterm}) in a covariant form:
\begin{align}
-2i m_f^2 A_\nu \sigma _\lambda \left( \hat{l}^\lambda \hat{l}^\nu \frac{1}{p^2}\ln{\frac{2p}{q}} + \hat{p}^\lambda \hat{p}^\nu \frac{1}{p^2} - \hat{p}^\lambda P^\nu \frac{1}{p^3}\ln{\frac{2p}{q}} \right), \label{vertex_without_thermal_mass}
\end{align}
with $\hat{p}^\nu = (0,\frac{\vec{p}}{p})$.%, $q^\nu = (\delta_\mu^\nu-n_\mu n^\nu) Q_\mu$.
 
As mentioned earlier, (\ref{vertex_without_thermal_mass}) still contains IR divergences as $q\to0$. To see the origin of the divergence, it is instructive to rewrite part of the integrand of \eqref{diag_abc} as\footnote{We could have used this representation directly in \eqref{diag_abc}, but this would make the origin of the divergence less clear.}
\begin{align}
\(\frac{1}{2K'\cdot P_2}-\frac{1}{2K'\cdot P_1}\)\frac{1}{2K'\cdot Q}=\frac{1}{2K'\cdot P_2}\frac{1}{2K'\cdot P_1}.
\end{align}
We see the cancellation of the factors $2K'\cdot Q$ in the numerator and denominator. This implies partial cancellation of the IR divergence among diagram (b) and (c). The remaining divergence arise in the kinematic region where $P_1$ and $P_2$ become collinear with $K'$ simultaneously. To divergence can be cut off by screening effect of the medium. The screening effect can be implemented by using dressed propagators. The discussion above suggest that we only need to dress the propagators with momenta $K'\pm P_1$ and $K'\pm P_2$. This amounts to the following substitutions \cite{le2000thermal}
%happens when $K \parallel P \&\& K \parallel P'$.
%So a resummation of self-energy for the internal propagators is needed. Again, it is not surprising, since in $q$ expansion, the fermion-background vertex can be regarded as a dressed propagator, and $K \parallel P_1 \&\& K \parallel P_2$ corresponds to the collinear divergence appeared in the one-loop self-energy. 
%\begin{figure}
%  \centering
    %\includegraphics[width=.8\textwidth]{figures/resummation.png}
%  \caption{the fermion-background vertex regarded as a dressed propagator. {\bf It is not clear to me what this means. Do we still need this Fig?}}
%\end{figure}
%Back to (\ref{sum_over_three}), we now introduce the resummation for divergent factor $\frac{1}{2K\cdot P'}$ and $\frac{1}{2K\cdot P}$. According to \cite{le2000thermal}, the resummation of fermion self-energy under HTL approximation gives:
\begin{align}
\frac{1}{2K'\cdot P_{1,2}} \to \frac{1}{2K'\cdot P_{1,2}+2m_f^2}, \label{substitute_1}
\end{align}
for fermions and
%and the resummation of gluon self-energy under HTL approximation gives\footnote{The change of polarization tensor does not change the Dirac structure\cite{Lin:2023ass}.}:
\begin{align}
\frac{-g_{\mu \nu }}{2K'\cdot P_{1,2}} \to \frac{\mathcal{P}^T_{\mu \nu }}{2K'\cdot P_{1,2}+m_g^2}, \label{substitute_2}
\end{align}
for gluons.\footnote{The longitudinal polarization does not contribute in the HTL approximation, leaving only the transverse polarization with $\mathcal{P}^T_{\mu \nu }$ being the transverse projector. The change of polarization tensor does not change the Dirac structure \cite{Lin:2023ass}.} Here $m_g^2=\frac{1}{3}g^2T^2(C_A+\frac{1}{2}N_f)$ is the gluon thermal mass, with $C_A=3$ and $N_f$ being the number of flavors in QGP.
 After applying (\ref{substitute_1}) and (\ref{substitute_2}) to the internal propagators in (\ref{diag_abc}), we obtain:
\begin{align}
\begin{split}
-2ig^2C_F A_\nu \int\frac{d^{4} K'}{( 2\pi )^{4}} 2K'^\nu (K' \cdot \sigma) \Big[ \frac{1}{2K'\cdot P_1+2m_f^2}\frac{1}{2K'\cdot P_2+2m_f^2}f(K'_0) \\
+ \left(\frac{1}{2K'\cdot P_2+m_g^2}-\frac{1}{2K'\cdot P_1+m_g^2}\right)\frac{1}{2K'\cdot Q}\tilde{f}(K'_0) \Big] 2\pi \delta (K'^2) ,
\end{split}
\end{align}
%note that we have not introduced thermal mass into the internal fermion factor  $\frac{1}{2K\cdot Q}$ of (b) and (c) in Figure \ref{vertex_correction}, since there is a degree of cancellation between the IR divergence in two diagrams. if we do introduce thermal mass to it, then the additional correction would only be at $\mathcal{O}(\frac{m^2}{k})\sim g^2T$ level, which is negligible compared to other corrections.
By repeating the previous procedure for integrations, we can finally arrive at:
\begin{align}\label{vertex_corr}
\delta \Gamma_{vertex} ^\nu A_\nu =& 2 m_f^2 A_\nu \sigma _\lambda \Bigg{[} \frac{1}{6 p^2}\(2 \ln \left(\frac{pT}{m_f^2}\right)+\ln \left(\frac{2 pT}{m_{g }^2}\right) 
-36 \ln (A)+\ln \left(16 \pi ^3 \right)+3\)\nonumber\\
&\times\left(\hat{l}^\lambda \hat{l}^\nu - \hat{p}^\lambda P^\nu \frac{1}{p} \right) + \hat{p}^\lambda \hat{p}^\nu \frac{1}{p^2} \Bigg{]},
\end{align}
with $A\simeq 1.282$ being the Glaisher constant. We can see the thermal masses of fermions and gluons effectively cut off the logarithmic divergences.

The remaining contribution to the EMFF is from the self-energy correction to the incoming/outgoing states. The corresponding diagrams are collected in Figure~\ref{self-energy}.
%A complete medium correction should also include the resummed self-energy correction. 
\begin{figure}\label{self-energy}
  \centering
    \includegraphics[width=.8\textwidth]{figures/self-energy_1.png}    
    \includegraphics[width=.8\textwidth]{figures/self-energy_2.png}
  \caption{self-energy corrections to both incoming and outgoing states, each with two different labelings in $ra$.}
\end{figure}
These self-energy corrections are actually similar to the vacuum counterpart. That is, the self-energy correction can be absorbed into the field strength renormalization. Under the condition $m_f^2 \ll p \ll T$, we have \cite{le2000thermal}:
\begin{align}\label{self_corr}
\delta \Gamma_{self-energy} ^\nu A_\nu = 2 \frac{m_f^2}{2p^2}\left( 1-ln\frac{2p^2}{m_f^2} \right) \sigma ^\nu A_\nu .
\end{align}

Now we can combine \eqref{vertex_corr} and \eqref{self_corr}, and use the decomposition \eqref{parameterization} to find the leading medium correction to the three form factors:
\begin{align}
\delta F_0 &= \frac{2m_f^2}{k^2}X+\frac{m_f^2}{k^2}\left( 1-ln\frac{2k^2}{m_f^2} \right), \label{delta_F_0}\\
\delta F_1 &= \frac{2m_f^2}{k^2}(X-1)+\frac{m_f^2}{k^2}\left( 1-ln\frac{2k^2}{m_f^2} \right), \label{delta_F_1}\\
\delta F_2 &= \frac{2m_f^2}{k^2}X+\frac{m_f^2}{k^2}\left( 1-ln\frac{2k^2}{m_f^2} \right), \label{delta_F_2}
\end{align}
where $$X=\frac{1}{6}\left(2 \ln \left(\frac{kT}{m_f^2}\right)+\ln \left(\frac{2 kT}{m_{g}^2}\right)-36 \ln (A)+\ln \left(16 \pi ^3 \right)+3\right),$$
and we have replaced $p$ with the kinetic momentum $k$ to account for effect of the external EMF at $\mathcal{O}(q^0)$.
%
Let us discuss the spin polarization effect based on \eqref{i_physical_meanings_of_form_factors}:\\
%
i. In vacuum all three form factors are degenerate as a consequence of vacuum interaction. In the medium, the degeneracy is expected to be lifted. Interestingly, our explicit results show partial lift with $\d F_1\ne\d F_2=\d F_0$, indicating the spin couplings to parallel and perpendicular magnetic field are different, but the spin couplings to perpendicular electric and parallel magnetic field remain the same. We believe the remaining degeneracy is a consequence of the HTL approximation, which may be fully lifted beyond the approximation.\\
%
ii. The medium corrections to all EMFF are real, which is consistent with our analysis based on time-reversal symmetry. This can be traced back to our requirement that the EMF does no work on the fermions, so that no dissipation is involved. Relaxing the requirement would necessarily lead to complex EMFF, but also with more complicated structures of the vertex.\\
%
iii. We have not considered self-energy correction to EOM. It is actually not needed: because chiral symmetry is unbroken and respected by the self-energy $\S\propto P\cdot \s$ \footnote{It is true for the particle excitation we consider, but not so for plasmino.} thus the self-energy corrected EOM is the same as the free EOM up to an overall normalization. This justifies the use of spinor solution to the free EOM as the incoming/outgoing states.
%From (\ref{delta_F_0}, \ref{delta_F_1}, \ref{delta_F_2}) we see that spin-magnetic coupling has a non-zero contribution along the direction $\hat{p}^i$. 



%\vskip 5mm


\section{Summary and outlook}\label{sec_outlook}

We have rederived the fermion correlation function in external electromagnetic field obtained in chiral kinetic theory using field theory by resumming tree level diagrams. The field theory method allows us to study radiative corrections to spin couplings. By imposing the condition that the electromagnetic field does no work to fermions, we find the correlation function can be equivalently interpreted as scatterings amplitude. Focusing on a thermal medium, we provide general decomposition of the correlation function using three electromagnetic form factors, corresponding to spin couplings to perpendicular electric field, perpendicular and parallel magnetic field respectively. We calculate QCD radiative corrections to the form factors in QGP under the HTL approximation as an example. We find the radiative correction leads to partial lift of the degeneracy of the form factors in vacuum. In particular, the spin couples to parallel and perpendicular magnetic field differently, while the spin couples to perpendicular electric and magnetic field with the same magnitude.

It would be interesting to relax the no-work restriction, to study spin coupling with more general field configurations. This will necessarily lead to more structures of interaction vertex and complex form factors, which will enable us to study spin polarization as well as spin relaxation phenomenon.

While the bulk of the analysis in this paper relies on perturbation theory, the interpretation of scattering and formulation using form factors is not restricted to weakly interacting systems and also not restricted to elementary particles. An interesting application is on composite particles. It is known that the vacuum electromagnetic form factors characterize the charge and spin distributions of the composite particle as seen by external electromagnetic fields. In-medium electromagnetic form factors encodes how the medium modifies the charge and spin distributions. We leave this for future studies.


%{\large {\bf Acknowledgment}}
\section*{Acknowledgments}
This work is in part supported by NSFC under Grant Nos 12075328, 11735007.

\appendix

\section{Electromagnetic field resummation at tree level }\label{sec_app_A}
In this appendix, we perform the detailed calculations of the Wigner function slightly rewritten in the following form
\begin{align}\label{Wigner}
S^{ra}(K) &= \int_s e^{iK \cdot s} S^{ra}(X,s)U(x_2,x_1) \nonumber\\
&= \int_s e^{iK \cdot s} \int_P e^{-iP \cdot s} S^{ra}_{re}(P)e^{iA \cdot s},
\end{align}
where $\int_s= \int d^4s$, $\int_P= \int \frac{d^4P}{(2\pi )^4}$, and $S^{ra}_{re}(P)$ is the correlation function resummed to all order in $A$ and up to first order in the momentum inflow $q$. We split the contributions into zeroth and first orders in $q$. The zeroth order contribution corresponds to resummation of constant $A$ field, which can be done conveniently as 
\begin{align}
S_{re}^{ra}(P)&=\frac{iP\cdot\bar{\s}}{P^2+i\h}+\frac{iP\cdot\bar{\s}}{P^2+i\h}(-iA\cdot\s)\frac{iP\cdot\bar{\s}}{P^2+i\h}+\cdots\nonumber\\
			&=\frac{1}{iP\cdot\s}+\frac{1}{iP\cdot\s}(-iA\cdot\s)\frac{1}{iP\cdot\s}+\cdots\nonumber\\
			&=\frac{1}{i(P-A)\cdot\s}=\frac{i(P-A)\cdot\bar{\s}}{(P-A)^2+i\h}.
\end{align}
Combining with the gauge link and performing the remaining integrals, we obtain
\begin{align}\label{O1}
S_{ra}(K)&=S_{re}^{ra}(K+A)=\frac{iK\cdot\bar{\s}}{K^2+i\h},\nonumber\\
	&\Rightarrow\; S_{ra}^\m(K)=\frac{iK^\m}{K^2+i\h}.
\end{align}

%From Figure. \ref{tree_level_resum}, we obtain:
The first order contribution resums diagrams with all constant $A$ field except one with momentum inflow. We illustrate diagrammatic contributions in Figure~\ref{split}.
\begin{figure}
	\centering
	\includegraphics[width=.8\textwidth]{figures/split.png}
	\caption{The diagrammatic contributions are split into $O(q^0)$ and $O(q)$ corresponding to second and third lines. The encircled crosses and empty circles denote EMF with momentum inflow and without. The bare circles denote the difference of the two symbols, i.e. $O(q)$ contribution from the EMF. The double lines correspond to resummed correlation function in constant $A$ field.} \label{split}
\end{figure}
Using the diagram representation in Figure \ref{split}, we find it is given by the convolution of two resummed correlation functions and a vertex expanded to $O(q)$. Note that the momentum inflow causes an $O(q)$ shift of momentum at the vertex. The momenta before and after the inflow can be found from the following Fourier factors
\begin{align}
\int_{s,y}\int_{P_1,P_2}e^{-i(P-A) \cdot s}e^{-iP\cdot (x_1-x_2)}e^{iP_1\cdot(x_1-y)+iP_2\cdot(z-x_2)}e^{iq\cdot y},
\end{align}
with the last two factors from the Fourier transforms of two resummed correlation functions and EMF respectively. Writing $iP_1\cdot(x_1-y)+iP_2\cdot(z-x_2)=i\frac{1}{2}(P_1+P_2)\cdot(x_1-x_2)-i(P_1-P_2)\cdot z$, we find the $s$ and $y$ integrals fix $P_1$ and $P_2$ as
\begin{gather}
\frac{1}{2}(P_1+P_2)=P-A=K,\quad P_1=P_2+q\nonumber\\
\Rightarrow\ P_1=K+\frac{q}{2},\quad P_2=K-\frac{q}{2}.
\end{gather}
The diagrammatic identity leads to the following representation
\begin{align}\label{O1Oq}
\frac{iP_1\cdot\bar{\s}}{P_1^2+i\h}(-i A\cdot\s)\frac{iP_2\cdot\bar{\s}}{P_2^2+i\h}.
\end{align}
We need to keep the $O(q)$ contribution of \eqref{O1Oq}. The explicit representation of \eqref{O1Oq} indicates the effect of momentum inflow lies entirely in the shift of the momenta before and after the vertex. We expand the numerator and denominator of \eqref{O1Oq} respectively up to $O(q)$
\begin{align}
\frac{1}{2}\tr[\s^\m P_1\cdot\bar{\s}A\cdot\s P_2\cdot\bar{\s}]&\simeq 2K^\m K\cdot A-A^\m K^2+i \e^{\m\n\r\s} K^\r q^\s A_\n,\nonumber\\
\frac{1}{P_1^2+i\h}\frac{1}{P_2^2+i\h}&\simeq \frac{1}{(K^2+i\h)^2}+O(q^2).
\end{align}
Note that the $O(q^0)$ contribution is already included in \eqref{O1}. We then easily obtain the remaining $O(q)$ contribution as
\begin{align}\label{Oq}
O(q):\;S_{ra}^\m=\frac{ \e^{\m\n\r\s} K^\r q^\s A_\n}{(K^2+i\h)^2}.
\end{align}

We will also need the case when the vacuum vertex $\s^\m$ is replaced by the medium modified one $\G^\m$. The modification to the resummation procedure above is minimal. At $O(q^0)$, the form factor $F_2$ is irrelevant. Using \eqref{spinor_products}, we easily find
\begin{align}
u^\dg(P_1)\G^\m u(P_2)A_\m=u^\dg(P_1)\s^\m u(P_2)A^F_\m+O(q),
\end{align}
with $A_F^\m=(F_0 A_0,F_1 \vec{A})$. The effect of form factors at $O(q^0)$ is a renormalization factor to the gauge fields. With gauge link modified accordingly, we arrive at the following resummed correlation function
\begin{align}
S_{ra}^\m(K)=S_{re}^{ra}(K+A_F)=\frac{iK\cdot\bar{\s}}{K^2+\i\h}.
\end{align}
Note that $K=P-A_F$ is the kinetic momentum with medium effect taken into account through $A\to A_F$. At $O(q)$, we have a similar representation based on the diagrammatic identity
\begin{align}
\frac{iP_1\cdot\bar{\s}}{P_1^2+i\h}(-i A\cdot\G)\frac{iP_2\cdot\bar{\s}}{P_2^2+i\h}.
\end{align}

%\begin{align}
%S^{ra}_{resum}(P) &= \frac{iP'\cdot \bar{\sigma}}{{P'}^2+i \eta} + \frac{iP'\cdot \bar{\sigma}}{{P'}^2+i \eta}(-iA\cdot \sigma)\frac{iP\cdot \bar{\sigma}}{P^2+i \eta} + \frac{iP'\cdot \bar{\sigma}}{{P'}^2+i \eta}(-iA\cdot \sigma)\frac{iP\cdot \bar{\sigma}}{P^2+i \eta}(-iA\cdot \sigma)\frac{iP\cdot \bar{\sigma}}{P^2+i \eta}+\cdots \nonumber\\
%&= \frac{iP'\cdot \bar{\sigma}}{{P'}^2+i \eta}\left[1+ (-iA\cdot \sigma)\frac{iP'\cdot \bar{\sigma}}{P'^2+i \eta} + \cdots \right]\left[ 1+ (-iA\cdot \sigma)\frac{iP\cdot \bar{\sigma}}{P^2+i \eta} + \cdots \right] \nonumber\\
%&= \frac{iP'\cdot \bar{\sigma}}{{P'}^2+i \eta} + \frac{iP'\cdot \bar{\sigma}}{{P'}^2+i \eta}\frac{1}{1-\frac{(A\cdot\sigma) (P'\cdot \bar{\sigma})}{{P'}^2+i \eta}}\frac{\frac{(A\cdot\sigma) (P\cdot \bar{\sigma})}{{P}^2+i \eta}}{1-\frac{(A\cdot\sigma) (P\cdot \bar{\sigma})}{{P}^2+i \eta}} \nonumber\\
%&= \frac{iP'\cdot \bar{\sigma}}{{P'}^2+i \eta} + \frac{iP'\cdot \bar{\sigma}}{((P'-A)\cdot\sigma)(P'\cdot \bar{\sigma })+i \eta}(-iA\cdot \sigma)\frac{iP\cdot \bar{\sigma}}{((P-A)\cdot\sigma)(P\cdot \bar{\sigma })+i \eta} \nonumber\\
%&= \frac{iP'\cdot \bar{\sigma}}{{P'}^2+i \eta} + \frac{i}{(P'-A)\cdot\sigma+i \eta}(-iA\cdot \sigma)\frac{i}{(P-A)\cdot\sigma+i \eta} \nonumber\\
%&\equiv  \frac{iP'\cdot \bar{\sigma}}{{P'}^2+i \eta} + \frac{iK'\cdot \bar{\sigma}}{{K'}^2+i \eta}(-iA\cdot \sigma)\frac{iK\cdot \bar{\sigma}}{K^2+i \eta}
%\end{align}
%where the first terms contribute to the free propagator, in which $P'$ is no different from $P$. The above procedure has repeatedly calculated the $\mathcal{O}(q^0)$ order terms. An appropriate calculation would be to drop the $\mathcal{O}(q^0)$ order terms, and to calculate them individually:
%\begin{align}
%S^{ra}_{(0)}(P) &= \frac{iP\cdot \bar{\sigma}}{{P}^2+i \eta}(-iA\cdot \sigma)\frac{iP\cdot \bar{\sigma}}{P^2+i \eta} + \frac{iP\cdot \bar{\sigma}}{{P}^2+i \eta}(-iA\cdot \sigma)\frac{iP\cdot \bar{\sigma}}{P^2+i \eta}(-iA\cdot \sigma)\frac{iP\cdot \bar{\sigma}}{P^2+i \eta}+\cdots \nonumber\\
%&= \frac{iP\cdot \bar{\sigma}}{{P}^2+i \eta}\left[ (-iA\cdot \sigma)\frac{iP\cdot \bar{\sigma}}{P^2+i \eta} + \cdots \right] \nonumber\\
%&= \frac{iP\cdot \bar{\sigma}}{{P}^2+i \eta}\frac{1}{1-\frac{(A\cdot\sigma) (P\cdot \bar{\sigma})}{{P}^2+i \eta}} -\frac{iP\cdot \bar{\sigma}}{{P}^2+i \eta} \nonumber\\
%&= \frac{iP\cdot \bar{\sigma}}{((P-A)\cdot\sigma)(P\cdot \bar{\sigma })+i \eta} -\frac{iP\cdot \bar{\sigma}}{{P}^2+i \eta} \nonumber\\
%&= \frac{i}{(P-A)\cdot\sigma+i \eta} -\frac{iP\cdot \bar{\sigma}}{{P}^2+i \eta} \nonumber\\
%&\equiv  \frac{iK\cdot \bar{\sigma}}{K^2+i \eta}-\frac{iP\cdot \bar{\sigma}}{{P}^2+i \eta} 
%\end{align}
%So we obtain:

%\begin{align}
%S^{ra}(K) &= \int_s e^{iK \cdot s} \int_P e^{-iP \cdot s} \left(S^{ra}_{resum}+S^{ra}_{(0)}\right)(P)e^{iA \cdot s} \nonumber\\
%&= \int_s e^{iK \cdot s} \int_P e^{-iP \cdot s}\left( \frac{iK\cdot \bar{\sigma}}{{K}^2+i \eta} + \frac{iK'\cdot \bar{\sigma}}{{K'}^2+i \eta}(-iA\cdot \sigma)\frac{iK\cdot \bar{\sigma}}{K^2+i \eta} \right)e^{iA \cdot s} \nonumber\\
%&= \frac{iK\cdot \bar{\sigma}}{{K}^2+i \eta} + \frac{iK'\cdot \bar{\sigma}}{{K'}^2+i \eta}(-iA\cdot \sigma)\frac{iK\cdot \bar{\sigma}}{K^2+i \eta}
%\end{align}
%Through the same approach one can derive $S^{ar}_{(re)}(K)$:
%
%\begin{align}
%S^{ar}(K) &= \frac{iK\cdot \bar{\sigma}}{{K}^2-i \eta} + \frac{iK'\cdot \bar{\sigma}}{{K'}^2-i \eta}(-iA\cdot \sigma)\frac{iK\cdot \bar{\sigma}}{K^2-i \eta}
%\end{align}


\section{KMS relation in external electromagnetic field}\label{sec_app_B}

In this appendix, we derive the KMS relation in the presence of a slow-varying EMF. We shall see up to $O(q)$ it is formally the same the usual KMS relation. From \eqref{Wigner}, we know the resummed correlation function and the counterpart without EMF are related by
\begin{align}
S^{ra}(K)=S_{re}^{ra}(P=K+A),
\end{align}
for constant $A$ field. The same relation actually holds for correlation functions with other operator orderings. To see this, we note that the effect of constant $A$ on correlation function can be easily deduced from the Lagrangian ${\cal L}=\bar{\psi}(i\pd_\m-A_\m)\s^\m\psi$. With $i\pd_\m$ identified as canonical momentum $P_\m$, we see constant $A$ simply changes canonical momentum into kinetic one $K=P-A$. Furthermore, the gauge link in the correlation function fixes $K$ as argument of Wigner function. Therefore the usual KMS relation in free theory holds for the resummed correlation function. In particular
\begin{align}
S_{re}^<(P)=-e^{-\b K_0}S_{re}^>(P)\Rightarrow S^<(K)=-e^{-\b K_0}S^>(K).
\end{align}

At $O(q)$, as reasoned in Appendix~\ref{sec_app_A}, we shall keep the $O(q)$ contribution of the expression below
%In this paper，we have used the KMS relation for the background-coupled propagators, which is actually practicable as proven below. According to Figure \ref{tree_level_resum}, the resummed lesser propagator at tree level is given by:
\begin{align}
S^<_{re}(K) &= -\int d^4P_1 \left[ S_0^<(P_1)(-iA\cdot \sigma )S_0^F(P_2)-S_0^{\bar{F}}(P_1)(-iA\cdot \sigma )S_0^<(P_2) \right] \delta ^4(P_2+q-P_1) \nonumber\\
&= -\left[S_0^<(P_1)(-iA\cdot \sigma )S_0^F(P_2)-S_0^{\bar{F}}(P_1)(-iA\cdot \sigma )S_0^<(P_2)\right]_{P_1\to P_2+q},
\end{align}
with the following resummed correlation function at $O(q^0)$
\begin{equation}
\begin{aligned}
S^F_0 &= K\cdot \bar{\sigma}\left[ \frac{i}{K^2+i \eta }-2\pi \tilde{f}(K_0)\delta (K^2) \right] ,\\
S^{\bar{F}}_0 &= K\cdot \bar{\sigma}\left[ \frac{-i}{K^2-i \eta }-2\pi \tilde{f}(K_0)\delta (K^2) \right] ,\\
S^<_0 &= -2\pi K\cdot \bar{\sigma}\tilde{f}(K_0)\delta (K^2) ,\\
S^>_0 &= -2\pi K\cdot \bar{\sigma}(1-\tilde{f}(K_0))\delta (K^2).
\end{aligned}
\end{equation}
%where we discuss only particles with positive energy. 
We first find a vanishing contribution from the second terms of $S^F$ and $S^{\bar{F}}$:
\begin{align}
\left[ 2\pi \tilde{f}(P_1^0)\delta (P_1^2)2\pi \tilde{f}(P_2^0)\delta (P_2^2)-2\pi \tilde{f}(P_1^0)\delta (P_1^2)2\pi \tilde{f}(P_2^0)\delta (P_2^2) \right](P_1 \cdot \sigma )(-iA \cdot \sigma )(P_2 \cdot \sigma ) = 0,
\end{align}
which means that in $S^F$ and $S^{\bar{F}}$ only the vacuum part contribute to the resummed correlation function, thus we have %In view of Schwinger-Keldysh contour, these vacuum terms do not have a weight factor along the imaginary axis, so in the thermal average of the correlator, they only contribute to a temperature-independent factor. the resummed propagator thus can be given as:
\begin{align}
S^<_{re}(K) &= -\left[S_0^<(P_1)(-iA\cdot \sigma )\frac{iP_2\cdot \bar{\sigma }}{P_2^2+i \eta }-\frac{-iP_1\cdot \bar{\sigma }}{P_1^2-i \eta }(-iA\cdot \sigma )S_0^<(P_2)\right]_{P_1\to P_2+q}\nonumber\\
&=-2\p\left[-\tilde{f}(P_1^0)\d(P_1^2)(-iA\cdot \sigma )\frac{iP_2\cdot \bar{\sigma }}{P_2^2+i \eta }-\frac{-iP_1\cdot \bar{\sigma }}{P_1^2-i \eta }(-iA\cdot \sigma )\tilde{f}(P_2^0)\d(P_2^2)\right]_{P_1\to P_2+q},
\end{align}
where we have used the KMS relation for free propagators in the second line.
We shall relate the above to $S_{re}^{ra}$ and $S_{re}^{ar}$, which are given by
\begin{align}
&S_{re}^{ra}=\frac{iP_1\cdot\bar{\s}}{P_1^2+i\h}(-iA\cdot\s)\frac{iP_2\cdot\bar{\s}}{P_2^2+i\h},\nonumber\\
&S_{re}^{ar}=\frac{iP_1\cdot\bar{\s}}{P_1^2-i\h}(-iA\cdot\s)\frac{iP_2\cdot\bar{\s}}{P_2^2-i\h}.
\end{align}
Using the following identity
\begin{align}
&\frac{1}{P_1^2+i\h}\frac{1}{P_2^2+i\h}-\frac{1}{P_1^2-i\h}\frac{1}{P_2^2-i\h}\nonumber\\
=&\(-2i\p\d(P_1^2)\)\frac{1}{P_2^2+i\h}-\frac{1}{P_1^2-i\h}\(2i\p\d(P_2^2)\),
\end{align}
we easily find the KMS relation holds if $\tilde{f}(P_1^0)=\tilde{f}(P_2^0)$. Recall that we are supposed to keep the $O(q)$ contributions, thus we may replace $P_1^0$ and $P_2^0$ by $K_0$ to obtain
%Using the KMS relation for free propagator: $S_0^<(K)=-e^{-\beta K_0}S_0^>(K)$, we have:
%\begin{align}
%S^<_{re}(K) =& \left[ e^{-\beta K'_0} S_0^>(K')(-iA\cdot \sigma )\frac{iK\cdot \bar{\sigma }}{K^2+i \eta }-\frac{-iK'\cdot \bar{\sigma }}{K'^2-i \eta }(-iA\cdot \sigma )e^{-\beta K_0}S_0^>(K)\right]_{K'\to K+q} \nonumber\\
%=& \bigg{[} e^{-\beta K'_0} S_0^>(K')(-iA\cdot \sigma )\left(P.V.\frac{iK\cdot \bar{\sigma }}{K^4 + \eta^2 } + \pi K\cdot \bar{\sigma }\delta (K^2) \right) \nonumber\\
%&+ \left(P.V.\frac{iK'\cdot \bar{\sigma }}{K'^4 + \eta^2 } - \pi K'\cdot \bar{\sigma }\delta (K'^2) \right)(-iA\cdot \sigma )e^{-\beta K_0}S_0^>(K)\bigg{]}_{K'\to K+q}
%\end{align}
%The above procedure is also applicable to $S^>_{re}$: 
%\begin{align}
%S^>_{re}(K) =& \left[S_0^>(K')(-iA\cdot \sigma )S_0^{\bar{F}}(K)-S_0^{F}(K')(-iA\cdot \sigma )S_0^>(K)\right]_{K'\to K+q} \nonumber\\
%=& \left[S_0^>(K')(-iA\cdot \sigma )\frac{-iK\cdot \bar{\sigma }}{K^2-i \eta }-\frac{iK'\cdot \bar{\sigma }}{K'^2+i \eta }(-iA\cdot \sigma )S_0^>(K)\right]_{K'\to K+q} \nonumber\\
%=& -\bigg{[} S_0^>(K')(-iA\cdot \sigma )\left(P.V.\frac{iK\cdot \bar{\sigma }}{K^4 + \eta^2 } - \pi K\cdot \bar{\sigma }\delta (K^2) \right) \nonumber\\
%&+ \left(P.V.\frac{iK'\cdot \bar{\sigma }}{K'^4 + \eta^2 } + \pi K'\cdot \bar{\sigma }\delta (K'^2) \right)(-iA\cdot \sigma )S_0^>(K)\bigg{]}_{K'\to K+q}
%\end{align}
%Using the condition $q_0=0$, we have $K'_0=K_0$ and $\tilde{f}(K'_0)=\tilde{f}(K_0)$, which will lead to the cancellation of the non-principle value part in the vacuum terms:
%\begin{align}
%S^<_{re}(K) =& e^{-\beta K_0} \bigg{[} -2\pi (1-\tilde{f}(K_0))(K \cdot \bar{\sigma })\delta (K'^2)(-iA\cdot \sigma )\left(P.V.\frac{iK\cdot \bar{\sigma }}{K^4 + \eta^2 } + \pi K\cdot \bar{\sigma }\delta (K^2) \right) \nonumber\\
%&+ \left(P.V.\frac{iK'\cdot \bar{\sigma }}{K'^4 + \eta^2 } - \pi K'\cdot \bar{\sigma }\delta (K'^2) \right)(-iA\cdot \sigma )(-2\pi (1-\tilde{f}(K_0))(K \cdot \bar{\sigma })\delta (K^2))\bigg{]}_{K'\to K+q} \nonumber\\
%=& e^{-\beta K_0}\left[ S_0^>(K')(-iA\cdot \sigma )P.V.\frac{iK\cdot \bar{\sigma }}{K^4 + \eta^2 } + P.V.\frac{iK'\cdot \bar{\sigma }}{K'^4 + \eta^2 }(-iA\cdot \sigma )S_0^>(K)\right]_{K'\to K+q}
%\end{align}
%\begin{align}
%S^>_{re}(K) =& -\bigg{[} -2\pi (1-\tilde{f}(K_0))(K \cdot \bar{\sigma })\delta (K'^2)(-iA\cdot \sigma )\left(P.V.\frac{iK\cdot \bar{\sigma }}{K^4 + \eta^2 } - \pi K\cdot \bar{\sigma }\delta (K^2) \right) \nonumber\\
%&+ \left(P.V.\frac{iK'\cdot \bar{\sigma }}{K'^4 + \eta^2 } + \pi K'\cdot \bar{\sigma }\delta (K'^2) \right)(-iA\cdot \sigma )(-2\pi (1-\tilde{f}(K_0))(K \cdot \bar{\sigma })\delta (K^2))\bigg{]}_{K'\to K+q} \nonumber\\
%=& -\left[ S_0^>(K')(-iA\cdot \sigma )P.V.\frac{iK\cdot \bar{\sigma }}{K^4 + \eta^2 } + P.V.\frac{iK'\cdot \bar{\sigma }}{K'^4 + \eta^2 }(-iA\cdot \sigma )S_0^>(K)\right]_{K'\to K+q}
%\end{align}
%Finally we obtain
\begin{align}
S_{re}^<(K)=-f(K_0)(S_{re}^{ra}(K)-S_{re}^{ar}(K)),
\end{align}
with a suppressed $O(q^2)$ correction to $S_{re}^<(K)$.

Finally let us briefly comment on the KMS relation when the vertex changes as $\s^\m\to\G^\m$. As we reasoned in Appendix~\ref{sec_app_A}, at $O(q^0)$, the effect is a simple renormalization of gauge field $A\to A_F$, which can be understood with the effective Lagrangian ${\cal L}=\bar{\psi}(i\pd_\m-A^F_\m)\s^\m\psi$. Based on the effective Lagrangian, the derivation in the above still applies to give formally the same KMS relation. At $O(q)$, the modification to the derivation is through $A\cdot\s\to A\cdot\G$. Note that the derivation in the above only uses the KMS relation at $O(q^0)$, not the detail of the vertex. So we again find the KMS relation holds.

%\begin{footnotesize}
%\begin{thebibliography}{999}%\refmark


%\end{thebibliography} %\end{footnotesize}
%\end{multicols}


\bibliographystyle{unsrt}\bibliography{EM_vertex}

\end{CJK}

\end{document}
